% ------------------------DIGITAL TWINS------------------

% Seminal
tuegel2011reengineering - Tuegel et al. (2011) introduced one of the first detailed visions of a Digital Twin (DT) as a tail-specific, ultra-high-fidelity structural model that evolves in response to both simulated and real flight data, enabling continuous aircraft life prediction. They proposed a fully integrated, multiphysics DT architecture combining CFD, FEM, thermal, stress, vibration, and damage models to simulate material degradation over time and update prognostics based on actual usage via Bayesian updating. The paper emphasized the shift from cycle-based to time-continuous, multiscale damage modeling, and outlined the major technical challenges in high-performance computing, uncertainty quantification, and data integration needed to operationalize DTs. It framed the DT not just as a predictive tool but as a unified decision-making system for maintenance, configuration control, and mission readiness in aerospace systems.

glaessgen2012digital - This paper defines the Digital Twin as an integrated, multiphysics, multiscale, probabilistic simulation of an as-built vehicle that continuously updates using sensor data, historical records, and fleet-wide information to mirror and predict the real-time behavior and degradation of its physical counterpart. The authors contrast this with conventional certification and sustainment practices rooted in heuristic safety factors, empirical design rules, and assumed similitude, which are inadequate for the extreme conditions expected in future aerospace missions. The Digital Twin paradigm supports virtual certification, in-situ forensics, and mission adaptability by enabling real-time health assessment, remaining useful life estimation, anomaly detection, and mission reconfiguration. Key contributions include a shift from reactive to proactive safety and performance assurance, integration with on-board IVHM systems, and alignment with national priorities such as Simulation-Based Engineering Science (SBES) and Integrated Computational Materials Engineering (ICME).

tao2018product - propose a comprehensive framework for digital twin (DT)-driven product design, manufacturing, and service, addressing major data fragmentation issues across the product lifecycle. The paper defines the digital twin as an integration of physical product, virtual model, and connected data, emphasizing real-time synchronization, interaction, and self-evolution. Key contributions include detailed DT-driven processes for design (conceptualization, detailed design, and virtual verification), manufacturing (through Digital Twin Shop Floors), and service (predictive maintenance, behavior analysis, and optimization strategies). The authors also identify five core technical challenges—ranging from intelligent sensing to virtual-real data mapping—and present case studies to validate the feasibility of the DT approach.

boschert2016digital - This chapter conceptualizes the Digital Twin (DT) as a dynamic, simulation-driven model architecture that mirrors the physical and functional characteristics of a product or system across all lifecycle phases—from design and engineering to operation and service. It defines the DT as a modular, evolving set of executable simulation models integrated with engineering data, operational inputs, and behavior representations, designed to support optimization, validation, and predictive maintenance. The authors emphasize the importance of model-based systems engineering (MBSE) and structured DT architecture in enabling consistent model reuse, system integration, and lifecycle-wide assistance functions such as condition monitoring and service planning. The DT is positioned not as a monolithic data store but as a selective, purpose-driven simulation toolchain that closes the loop between physical systems and virtual models, facilitating faster time-to-market, better quality, and smarter service strategies in complex mechatronic and cybertronic systems.

rosen2015importance - This paper defines the Digital Twin as a structured, simulation-driven digital replica of physical assets and systems that evolves across the lifecycle and serves as a core enabler of autonomy in manufacturing. It positions modularity, connectivity, autonomy, and digital twins as the four foundational pillars of Industrie 4.0, emphasizing that autonomous systems require detailed digital models to perform intelligent, decentralized decision-making in dynamic environments. Through examples in cyber-physical production systems (CPPS), the paper demonstrates how digital twins facilitate self-adaptive behavior, such as decentralized scheduling, real-time decision support, and fault tolerance. The Digital Twin is shown to bridge lifecycle stages by unifying design-time simulation with operational data, enabling predictive maintenance, flexible reconfiguration, and continuous system optimization.


% Definitions, Elaborations
tao2018digital - This paper provides a comprehensive review of digital twin (DT) research and industrial applications, covering 50 academic articles, 8 patents, and multiple industry case studies. It identifies Prognostics and Health Management (PHM) as the most developed and promising area, where DTs outperform traditional methods through real-time data fusion, high-fidelity modeling, and interactive virtual-physical feedback. The authors emphasize that DT modeling is the core technical challenge, noting the lack of a unified, standardized modeling framework that incorporates all five essential DT dimensions: physical, virtual, data, service, and connection. Another critical challenge is cyber–physical fusion, which requires advancements in data acquisition, real-time analytics, and secure communication protocols. The paper concludes by recommending future research into DT-driven optimization, operational control, and fault-tolerant, scalable architectures for smart manufacturing systems.

kritzinger2018digital - conduct a structured literature review of Digital Twin (DT) applications in manufacturing and introduce a three-level taxonomy—Digital Model (DM), Digital Shadow (DS), and Digital Twin (DT)—based on the degree of integration between physical and digital systems. They highlight that while many papers claim to use DTs, most implementations actually fall into the lower-integration categories of DM or DS, with true bidirectional DTs being rare. The review identifies Production Planning and Control (PPC) and maintenance as the most common application areas and emphasizes that most contributions are conceptual rather than practical. The authors call for clearer definitions, more standardized reference models, and an increase in real-world case studies to advance DT maturity in manufacturing.

grieves2017digital - This foundational paper defines the Digital Twin (DT) as a dynamic virtual representation of a physical system that evolves across the entire product lifecycle, aimed at identifying and mitigating unpredicted undesirable (UU) emergent behaviors in complex systems. It introduces a taxonomy of emergent behavior based on Holland’s classification and distinguishes between static and evolutionary emergence, highlighting that many catastrophic failures stem from latent design-time issues or unpredictable human interaction. Key contributions include a four-quadrant framework of system behavior (PD, PU, UD, UU), a lifecycle-based DT model (with DTP, DTI, DTE constructs), and the Grieves Tests of Virtuality (Visual, Performance, Reflectivity) to assess DT fidelity. The paper argues that DTs, supported by simulation, real-time feedback, and system "front-running," can dramatically reduce lifecycle waste and improve resilience in socio-technical systems like NASA’s space infrastructure.


semeraro2021digital - This paper defines a Digital Twin (DT) as a set of adaptive models that emulate the behavior of a physical system through real-time data, enabling simulation, prediction, and optimization across the product life cycle. It presents a structured paradigm based on a systematic review of 150 papers, organizing DT research along five dimensions: application contexts, life cycle phases, functions, architecture layers, and enabling components. Using text mining and Formal Concept Analysis, the study identifies common DT architectures (physical, network, computing layers) and major application domains (manufacturing, healthcare, aerospace, city management, maritime). The review highlights open challenges including the need for modular, interoperable architectures and broader support for lifecycle-wide DT integration.


tomczyk2022digital - This paper analyzes the evolution of Digital Twin (DT) definitions through a structured literature review and introduces a concept matrix to compare conceptual streams. It identifies a paradigm shift from the classic three-dimensional DT model—physical space, virtual space, and a bidirectional data link (Grieves, Glaessgen \& Stargel)—to an extended five-dimensional model that adds data and services as core dimensions (Tao et al.). The authors also distinguish an “applicational” approach, where DT definitions are shaped by use-case context. The study contributes a taxonomy of DT definitions, highlights the growing centrality of data, and shows that recent literature increasingly aligns with the extended model, emphasizing DT’s dynamic and scalable nature across lifecycles and domains.

gabor2016simulation - This paper proposes a simulation-based architecture for smart cyber-physical systems (SCPS) that tightly integrates digital twins—ultra-high-fidelity simulations capable of mimicking system behavior across space and time scales—for real-time planning, optimization, and fault detection. The authors define a modular software framework in which digital twins share the same interfaces as physical components, allowing seamless interchangeability between simulated and real environments. A key contribution is the tiered control model (Tiers 0–3) that classifies system behavior constraints from physical laws to planning-level objectives, enabling systematic safety and autonomy engineering. The architecture also supports multiple parallel digital twins and introduces the need for runtime learning and trust-based reconciliation of diverging simulation models.


% applications
liu2021review - This paper presents a comprehensive review of 240 digital twin (DT) publications, classifying them by lifecycle phase, integration level (Digital Model, Shadow, Twin), and application domain. It synthesizes the evolution of DT definitions—culminating in a view of DTs as individualized, high-fidelity, real-time, and bidirectionally controllable digital replicas—and identifies three enabling technological pillars: data acquisition and communication, high-fidelity multi-physics modeling, and model-based simulation. The authors categorize 15 types of industrial applications across design, manufacturing, service, and retire phases, emphasizing that most real-world implementations are limited to digital shadows, with full lifecycle and bidirectional integration still rare. They highlight key research gaps, such as insufficient full-lifecycle integration, limited modeling of complex phenomena, and lack of unified frameworks for DT deployment across heterogeneous systems.


botin2022digital - This systematic review defines Digital Twins (DTs) as virtual replicas of physical systems with bidirectional data exchange, and categorizes them by type (Instance, Prototype, Performance), integration level (Model, Shadow, Twin), and maturity (from static models to autonomous systems). The authors analyze DT applications across five domains—smart cities, logistics, medicine, engineering, and automotive—highlighting enabling technologies like IoT, AI/ML, and remote sensing, and evaluate implementations using Technology Readiness Level (TRL), Societal Readiness Level (SRL), and maturity indices. Key challenges include the lack of standardization, data and communication complexity, high implementation costs, and low societal integration. The paper calls for cross-domain knowledge transfer, better stakeholder inclusion, and future research into scalable, interoperable, and socially responsible DT systems.

barricelli2019survey - This survey synthesizes 75 works to clarify Digital Twin (DT) definitions, key characteristics, application domains, and sociotechnical design implications. DTs are defined as intelligent, continuously evolving virtual counterparts of physical systems, distinguished by real-time data exchange, AI integration, and closed-loop optimization. The authors identify three core domains of DT deployment—manufacturing, aviation, and healthcare—while outlining a lifecycle-based design taxonomy that includes proactive and retroactive DT creation. A major contribution is the call for sociotechnical design practices and end-user development tools to support stakeholder collaboration, alongside identification of persistent challenges such as data ethics, interoperability, and unequal access to DT infrastructure.

javaid2023digital - This paper defines a Digital Twin (DT) as a real-time, virtual replica of physical systems, processes, or assets, integrated with IoT, AI, and data analytics to support monitoring, simulation, and optimization throughout the asset lifecycle. It proposes a structured framework of DT features, discusses the technological requirements (e.g., sensors, cloud platforms, canonical data models), and outlines its supportive role in Industry 4.0 for predictive maintenance, process optimization, and smart manufacturing. The authors identify 33 major application areas of DT in Industry 4.0 across domains such as healthcare, automotive, logistics, and smart cities, highlighting its ability to simulate reality, enable automation, reduce downtime, and support design and innovation. The paper concludes by emphasizing the transformative potential of DT in achieving scalable, cost-efficient, and intelligent operations, while also noting challenges related to data security, integration, and infrastructure investment.

